\begin{proof}
\textbf{1. Preliminaries and notation.}  
Set  
\[
\lambda:=\frac1p\in(0,1),\qquad 1-\lambda=\frac1q .
\]
Hence \(p=\frac1\lambda\) and \(q=\frac1{1-\lambda}\); the relation \(\frac1p+\frac1q=1\) is equivalent to \(q=\frac{p}{p-1}\) (or \(p=\frac{q}{q-1}\)).

\medskip\noindent\textbf{2. Degenerate cases.}  
If \(a=0\) or \(b=0\) then \(ab=0\) while the right–hand side of  
\[
ab\le\frac{a^{p}}{p}+\frac{b^{q}}{q}
\tag{1}
\]
is a sum of non‑negative terms, so (1) holds. Equality forces both terms on the right to vanish, i.e. \(a=b=0\).

\medskip\noindent\textbf{3. Convex function and its conjugate.}  
Consider  
\[
\varphi:[0,\infty)\to[0,\infty),\qquad \varphi(t)=\frac{t^{p}}{p}.
\]

\textit{Convexity.}  
\[
\varphi'(t)=t^{p-1},\qquad 
\varphi''(t)=(p-1)t^{p-2}\ge0\quad(p>1,\;t\ge0),
\]
so \(\varphi\) is convex.

\textit{Legendre (convex) conjugate.}  
\[
\varphi^{*}(s)=\sup_{t\ge0}\bigl\{\,st-\varphi(t)\,\bigr\},
\qquad s\ge0 .
\]
To locate the supremum differentiate the bracketed expression:
\[
\frac{d}{dt}\bigl(st-\tfrac{t^{p}}{p}\bigr)=s-t^{p-1}.
\]
The unique maximiser satisfies \(s=t^{p-1}\), i.e.  
\[
t=s^{\,q-1}\qquad\bigl(q-1=\tfrac1{p-1}\bigr).
\]
Substituting this \(t\) gives
\[
\begin{aligned}
\varphi^{*}(s)
&=s\cdot s^{\,q-1}-\frac{(s^{\,q-1})^{p}}{p} \\
&=s^{q}-\frac{s^{q}}{p}
   =\frac{s^{q}}{q},
\end{aligned}
\]
where we used \(\frac1p+\frac1q=1\).

\medskip\noindent\textbf{4. Fenchel–Young inequality.}  
For any convex \(\varphi\) and \(s,t\ge0\) one has  
\[
st\le\varphi(t)+\varphi^{*}(s).
\tag{2}
\]
(See Rockafellar, *Convex Analysis*, Theorem 23.5.)

Applying the concrete \(\varphi\) and \(\varphi^{*}\) from step 3 to (2) yields
\[
st\le\frac{t^{p}}{p}+\frac{s^{q}}{q}.
\tag{3}
\]

\medskip\noindent\textbf{5. Application to the present variables.}  
Choose \(t=a\) and \(s=b\) in (3). Then (3) becomes exactly (1):
\[
ab\le\frac{a^{p}}{p}+\frac{b^{q}}{q}.
\]

\medskip\noindent\textbf{6. Equality condition.}  
Equality in (2) holds iff the maximiser satisfies \(s\in\partial\varphi(t)\).
Since \(\varphi\) is differentiable on \((0,\infty)\),
\[
s=\varphi'(t)=t^{p-1}.
\]
With \(s=b\) and \(t=a\) this reads \(b=a^{p-1}\), i.e.  
\[
a^{p}=b^{q}.
\tag{4}
\]
Together with the degenerate case of step 2 we obtain the complete equality
description:
\[
a=b=0\quad\text{or}\quad a,b>0\ \text{and}\ a^{p}=b^{q}.
\]

\medskip\noindent\textbf{7. Alternative elementary proof (weighted AM–GM).}  
Write
\[
ab=(a^{p})^{1/p}\,(b^{q})^{1/q}.
\]
Since \(\frac1p+\frac1q=1\), the weighted AM–GM inequality with weights
\(\frac1p\) and \(\frac1q\) gives
\[
(a^{p})^{1/p}\,(b^{q})^{1/q}
\le\frac1p\,a^{p}+\frac1q\,b^{q},
\]
which is precisely (1). Equality in weighted AM–GM occurs exactly when the
two weighted terms are equal, i.e. when \((a^{p})^{1/p}=(b^{q})^{1/q}\), again
equivalent to (4).

Thus Young’s inequality holds for all \(a,b\ge0\) with conjugate exponents
\(p,q>1\), and equality occurs precisely in the cases described above. \(\square\)
\end{proof}